\documentclass[titlepage]{article}
\usepackage{graphicx,wrapfig} % Required for inserting images and arranging them in text
\usepackage{amsmath} % For matrices and other maths stuff, very good
\usepackage{amsthm} % For \newtheorem stuff. Makes nice definitions and theorems
\usepackage{amssymb} % For does not exists \nexists
\usepackage{multicol} % For side - by - side equations/text
\usepackage{amsfonts} %For set notation, like R
\usepackage{xfrac} %fraction sideways
\usepackage{hyperref} %Clickable table of contents and references and hyperlinks
\usepackage{tikz, pgfplots} %Nice diagrams and graphics
\usepackage{mdframed} %Boxes around theorems
\usepackage{fancyhdr} %Fancy Header and Footer
\usepackage{tensor} %For tensor indices. Used mainly with \indices{}
\usepackage{physics} %For dirac notation
\usepackage{mathtools} %Text above \rightarrow
\usepackage{enumitem} %Editable enumerate items
\usepackage{xparse} %For Definition
%\usepackage{newpxmath} %For improved spacing between math symbols
\usepackage{bm} %Shorthand for bold text in mathmode


\usepackage[T1]{fontenc}
\usepackage[utf8]{inputenc}
\usepackage[lithuanian]{babel}

\usepackage[compat=1.1.0]{tikz-feynman} %For Feynman diagrams
\usetikzlibrary{positioning}

\usepackage[a4paper, total={6in, 8in}]{geometry} %Formating margins
\pagestyle{fancy}
\setlength{\headheight}{15pt}
\fancyhead[l]{Gytis Vejelis}
\fancyhead[r]{}
\fancyhead[c]{\textbf{Baryogenesis}}
\fancyfoot[c]{\thepage}

\NewDocumentCommand{\Definition}{ O{Definition} m m }{%
  \par\noindent\textbf{#1 #2.}\;#3\par
}

\newenvironment{definition}[1][Definition]{%
  \vspace{1em}%
  \begin{defn}[#1]%
}{%
  \end{defn}%
  \vspace{1em}%
}

\def\Include#1{%
        \def\ChapterPath{#1/}%
        \def\GraphicsPath{\ChapterPath Images/}%
        \include{\ChapterPath#1}}

\newcommand\IncludeGraphics[2][width=\linewidth]{%
    \includegraphics[#1]{\GraphicsPath #2}} 

\newcommand\bld[1]{\textbf{#1}}
\newcommand\basis[1]{\text{$\hat{\textbf{#1}}$}}
\newcommand\set[1]{\mathbb{#1}}
\newcommand\setn[2]{\mathbb{#1}^{#2}}
\newcommand\mpl{M_{\rm pl}}

\newcommand\de[2]{\frac{d#1}{d#2}}

\hypersetup{
    colorlinks,
    citecolor=black,
    filecolor=black,
    linkcolor=black,
    urlcolor=black
}


\title{Atom Theory}
\author{Gytis Vejelis}
\date{2026}

\pgfplotsset{compat=1.18}
\begin{document}

\maketitle

\tableofcontents

\newpage

\section{General Info}

Egzaminas - bilietai, 2 klausimai, 1 valanda + aptarimas. 

Paskutinė paskaita - klausimų aptarimas.

Galima daryti calculation su programa/pratybų viduryje semestro, tada egzamine bus tik vienas klausimas.

\section{1 Paskaita}

\subsection*{Pagrindinės atomų savybės}

Atomas neturi apibrėžtos ribos, nes jo sandarą lemianti elektrostatinė jėga yra toliasiekė

\begin{itemize}
    \item Mažiausias matomas dalykas $10^{-6}m$\
    \item Atomas $10^{-6}m$
    \item Branduolys $10^{-15}m$
\end{itemize}

Atomo matmenys apribūdina:
\begin{itemize}
    \item Boro orbitos spindulys $a$
    \item Vidutinis išorinio elektrono atstumas nuo branduolio
\end{itemize}

Vandenilio atome:
\begin{itemize}
    \item Pirmosios Boro orbitos spindulys
        \begin{equation*}
            a_0 = \frac{\hbar}{m_e e^2} \simeq 0.5 \cdot 10^{-10}m
        \end{equation*}
    \item $n$-tosios Boro orbitos spindulys
        \begin{equation*}
            a_{0n} = n^2 a_0 
        \end{equation*}
\end{itemize}

Vidutinis elektrono atstumas nuo branduolio:

\begin{equation*}
    \bar{r} = \int \Psi*(\bm{r}) \bm{r} \Psi{\bm{r}} d \bm{r} \qquad \bar{r}_{nl} ~ \frac{1}{Z} 
\end{equation*}
čia $nl - 1s,2s,2p,3s,3p,3d,4s$. $n$ - pagrindinis kvantinis skaičius (1,2,3,...), $l$ - orbitinis kvantinis skaičius (s,p,d,...)

Atomo masė: $1u \simeq 1.66057 \cdot 10^{-27} kg$. Atominis masės vienetas - tai anglies izotopo ${}^{12}C$ atomo massės $\frac{1}{12}$ dalis.

Atomo branduolio dydis: $R = (1.3 \pm 0.1)A^{\frac{1}{3}}10^{-15}m$, čia $A$ - masės skaičius (sveikas skaičius, artimiausias santykinei atomo masei.) Pvz
\begin{itemize}
    \item $R({}^{12}C)=2.7 \cdot 10^{-15} m$
    \item $R({}^{208}Pb)=7.10 \cdot 10^{-15} m$
\end{itemize}

Teigiamas jonas - atomas, netekęs elektronų
Neigiamas jonas - atomas, prisijungęs elektronus
Egzotinis atomas - prisijungęs pozitroną, miuoną arba kitokią elementariąją dalelę

Atomas, netekęs $n$ elektronų - $n$-kartinis jonas $A^{n+}$. Pvz $Ar^{4+} \equiv ArV$ - keturkartinis jonas.

Miuoniniai atomai - gaunami dirbtiniu būdu: protonais bomborduojant medžiagą; kai kurie sukurti miuonai yra pagribiami atomų. Miuonas - sunkusis elektronas, toks pat krūvis, 207 kartus didesnė masė. Miuonas nestabilus, suyra į kitas daleles per $10^{-6}s$. Dėl didelės masės miuonas sukasi labai arti branduolio. Jo spinduliutė labai priklauso nuo branduolio savybių - krūvio pasiskirstymo jame, formos, magnetinio momento. Miuoninių atomų spektrai teikia svarbią informaciją apie atomų branduolius.

Jonų susidarymas gamtoje ir laboratorijose. Žvaigždžių spektruose stebimos nuo kelių iki keliolikos jonizuotų atomų spektro linijos. Laboratorijose galima gauti bet kurio elemento, bet kurį joną.

EBIT - electron beam ion trap.

Atomas ir Saulės sistema:
\begin{itemize}
    \item Panašumai:
    \begin{itemize}
        \item Mažas masyvus centras
        \item Aplink skrieja mažesnės masės kūnai
    \end{itemize}
    \item Skirtumai:
    \begin{itemize}
        \item Atomas - sfera, Saulės sistema - diskas,
        \item stiprios elektrinės jėgos ir silna gravitacijos jėga
        \item elektronai identiški, planetos skirtingos
        \item Saulės sistemoje veikia klasikinė mechanika, tikslio padėtys, o atome - kvantinė mechanika, tikimybinis pasiskirstymas.
    \end{itemize}
\end{itemize}

Pagrindiniai dydžiai:
\begin{itemize}
    \item $a_0 = 0.52918 \cdot 10^{-10} m$ - pirmosios Boro orbitos spindulys
    \item $m_e=9.1095 \cdot 10^{-10} kg$
    \item $\hbar = \frac{h}{2 \pi} = 1.0546 \cdot 10^{-34} Js$ - veikimo vienetas
    \item $e=1.6022 \cdot 10^{-19} C$
\end{itemize}

Išvestiniai dydžiai:
\begin{itemize}
    \item Energijos vienetas 
    $\frac{m_e e^4}{\hbar^2} = 4.3438 \cdot 10^{-18} J$
    
    $1eV = 1.6022 \cdot 10^{-19}J$

    $1 a.v. = 27.213 eV$

    \item Laiko vienetas:
    $\frac{\hbar^3}{m_e e^4} = 2.4192 \cdot 10^{-16} s$

    \item Atominis greičio vienetas
    $\frac{e^2}{\hbar} = 2.1877 \cdot 10^6 m/s$

    $c \simeq 137 a.v.$
\end{itemize}


Lygtys:
\begin{itemize}
    \item Klasikinėje mechanikoje - II Niutono dėsnis
    \item Kvantinė mechanioje - Schrodingerio lygtis
\end{itemize}

Stacionarioji (nekintanti laike) Schrodingerio lygtis:

\begin{equation*}
    \hat{H} \Psi = E \Psi
\end{equation*}
čia $\hat{H}$ yra energijos operatorius Hamiltonianas, $E$ energija, $\Psi$ sistemos tikrin4 banginė funckija. Vieno laisvo atomo atveju paprastai priimamos prielaidos:
\begin{itemize}
    \item Atomo branduolys nejuda (su juo susijama koordinačių sistemos pradžia)
    \item Branduolys yra taškinis
\end{itemize}

%-----------------

Lygtis radialijai funkcijai gaunama, įstačius į vienelektronę Schrodingerio lygtį:

\begin{equation*}
    \hat{h} \psi_{nlm} = \epsilon_{nl} \psi_{nlm}
\end{equation*}

Hamiltoniano išraišką
\begin{equation*}
    \hat{h} ž \frac{p^2}{2m_e}
\end{equation*}


%------------------

Pertvarkius gaunasi tokia lygtis radialiajai funcijai (atominiais vientetais):
\begin{equation*}
    \left( \frac{d^2}{dr^2} - 2 V_{nl}(r) - \frac{l (l+1)}{r^2} \right) P_{nl}(r) = -2 \epsilon_{nl} P_{nl}(r)
\end{equation*}

Sprendinys turi tenkinti kraštines sąlygas: $P_{nl} (0) = 0$, $P_{nl}(\infty) = 0$.

Sprendiniai egzistuoja tik tam tikroms tikrinėms $\epsilon_{nl}$ vertėms, sudarančioms diskretinį sprektrą. Skaičius $n$ numeruoja lygties sprendinius, pradedant nuo $n=l+1$.

%----------------




\section{Vienelektronės Būsenos}

Sukinys ir sukinė funkcija

Elektronas turi dar sabajį judėjimo kiekio momentą - sukinį, kuris neturi analogo klasikinėje fizikoje. Nereliatyvistinėje teorijoje sukinio egzistavimą yra postuluojamas. Sukinio kvadrato $\hat{\bm{s}}^2$ ir šio momento projekcijos į $z$ ašį $\hat{s}_z$ operatoriai tenkina tikrinių verčių lygtis, analogiškai kaip $\hat{\bm{l}}^2$, $\hat{l}_z$:

\begin{align*}
    \hat{\bm{s}}^2 \chi_{\mu} (\sigma) = \hbar^2 s(s+1) \chi_{\mu} (\sigma)\\
    \hat{s}_z \chi_{\mu} (\sigma) = \hbar \mu \chi_{\mu} (\sigma)
\end{align*}
čia $\chi_{\mu}(\sigma)$ -šių operatorių tikrinė funkcija, $s$- sukinio kvantinis skaičius, lygus 1/2, o $\mu$ - sukinio projekcijos kvantinis skaičius, įgyjantis dvi vertes 1/2 ir -1/2
...

Nereliatyvistinis vienelektronis hamiltonianas $\hbar$ neveikia 5 sukinė kintamajį, tad $\hat{h}$, $\hat{\bm{s}}^2$ ir $\hat{s}_z$ komutuoja tarpusavyje, ir $\chi_{\mu}(\sigma)$ yra tikrinė $\hat{h}$ funckija. Taigi, atsižvelgiant į sukinį, vienelektronė nereliatyvistinė banginė funkcija užrašoma;

\begin{equation*}
    \psi_{}
\end{equation*}
...

Lyginumo transformacija - erdvinių koordinačių ženklo pakeitimas

\begin{equation*}
    x \rightarrow -x, \quad y \rightarrow -y, \quad z \rightarrow -z
\end{equation*}
Hamiltonianas $\hat{h}$ nesikei2ia, tad lyginumo operacij1 atitinkantis operatorius $\hat{P}$ komutuoja su hamiltonianu:

\begin{equation*}
    [\hat{P}, \hat{h}] = 0.
\end{equation*}

Vadinasi, hamiltoniano tikrin4 funkcija kartu yra ir $\hat{P}$ operatoriaus tikrin4 funkcija

\begin{equation*}
    \hat{P} \psi = c \psi
\end{equation*}
čia $c$ - lyginumo tikrinė vertė, kuri lygi -1 ir +1, nes $\hat{P}^2 \psi = c^2 \phi = \psi$

Banginės
...

Tikimybė rasti elektroną pasiskirstymas
\begin{equation*}
    \rho_{n l m \mu}(r, \theta, \phi, \sigma) = |\psi_{n l m \mu} (r, \theta, \phi, \sigma)|^2 = |\Theta_{lm}(\theta)|^2 \frac{P_{nl}^2 (r)}{r^2} \chi_{\mu}^2 (\sigma)
\end{equation*}

Atsižvelgiant, kad $Y_{lm} (\theta \phi) = \Theta_{lm} (\theta) e^{im \phi}$


\section(Daugiaelektronių atomų būsenos)

\subsection{Dvielektronio atomo hamiltonianas ir judėjimo kiekio momentai}

Atomo su dviem elektronais nereliatyvistinis hamiltonianas

\begin{equation*}
    \hat{H} = \frac{\hat{p}_1^2}{2m} +\frac{\hat{p}_2^2}{2m} - \frac{Z e^2}{4 \pi \epsilon_0 r_1} -\frac{Z e^2}{4 \pi \epsilon_0 r_2} + \frac{e^2}{4 \pi \epsilon_0 r_{12}}
\end{equation*}
čia $r_{12} = |\bm{r}_1 - \bm{r}_2 = \sqrt{(x_1 - x_2)^2 + (y_1 - y_2)^2 + (z_1 - z_2)^2} = \sqrt{r_1^2 + r_2^2 - 2 r_1 r_2 \cos \omega}$ 

...

Komutuojančių operatorių rinkinys. Hamiltonianas $\hat{H}$ neveikia 5 sukinius kintamuosius, tas jis komutuoja ne tik su $\hat{\bm{J}}^2$, $\hat{J}_z$, bet ir su $\hat{L}^2, \hat{L}_z$. Be to, $\hat{L}^2$ ir $\hat{L}_z$ komutuoja su atskirų eletroknų operatoriais $\hat{l}_i^2$, bet nekomutuoja su projekcijų operatoriais $\hat{l}_{iz}$:

Vadinasi, $\hat{\bm{L}}$, $\hat{L}_z$, $\hat{l}_1^2$, $\hat{l}_2^2$ sudaro komutuojančių operatorių rinkinį. Arsižvelgiant į elektronų sukinius, prasideda $\hat{\bm{S}}^2$, $\hat{S}_z$, $\hat{\bm{s}}_2$

...

Ji sudaroma iš atskirų elektronų banginių funkcijų. Jų sandaugos

\begin{equation*}
    \psi (l_1 m_1 | \bm{r}_1) \psi(l_2 m_2 | \bm{r}_2)
\end{equation*}
nėra tikrinės operatorių $\hat{\bm{L}}^2$ it $\hat{L}_z$ funkcijos, bet pastarosios gaunamos, imant toki7 sandaug7 tiesines kombinacijas:

\begin{equation*}
    \Psi(l_1, l_2 LM|\bm{r}_1 \bm{r}_2) = \sum_{m_1, m_2} 
    \begin{bmatrix*}
        l_1 & l_2 & L\\
        m_1 & m_2 & M
    \end{bmatrix*}
    \psi (l_1, m_1 | \bm{r}_1) \psi (l_2, m_2 | \bm{r}_2)
\end{equation*}

Skleidimo koeficientai yra vadinami Clebsch-Gordano (Klebšo ir Gordano) koeficientais.
Atsižvelgiant ir į sukinius:

\begin{equation*}
    \Psi(l_1, s_1, l_2 s_2 L S M_L M_S |\tau_1 \tau_2) = \sum_{m_1, m_2, \mu_1, \mu_2} 
    \begin{bmatrix*}
        l_1 & l_2 & L\\
        m_1 & m_2 & M_L
    \end{bmatrix*}
    \begin{bmatrix*}
        s_1 & s_2 & S\\
        m_1 & m_2 & M_S
    \end{bmatrix*}
    \psi (l_1 s_1, m_1 \mu_1 | \tau_1 ) \psi (l_2 s_2, m_2 \mu_2 | \tau_2 )
\end{equation*}

\subsubsection*{Clebsch-Gordano koeficientas}

\begin{equation*}
    \begin{bmatrix*}[c]
        j_1 & j_2 & j \\
        m_1 & m_2 & m
    \end{bmatrix*}
\end{equation*}

Savybės:

$j_i$, $m_i$, taip pat $j$, $m_i$ turi b8ti abu sveiki abu pusiniai skaičiai.

\begin{itemize}
    \item $|j_1 - j_2| \leq j \leq j_1 + j_2$
    \item $m_i \leq j_i$
    \item $m = m_1 + m_2$
\end{itemize}

\begin{equation*}
  \begin{bmatrix*}[c]
    j_1 & j_2 & j \\
    m_1 & m_2 & m
  \end{bmatrix*}
  = \delta_{m_1+m_2,\,m}\,\Delta(j_1 j_2 j)\,
  \left[\frac{(j+m)!\,(j-m)!\,(2j+4)!}{\cdots}\right]
\end{equation*}

...

Surištų momentų funkcija $\Psi(l_1 s_1 l_2 s_2 L S M_L M_S)$ netenkina:
\begin{itemize}
    \item Paulio principo: du elektronai negali būsenose su visais tais pačiais kvantiniais skaičiais;
    \item tapatingumo principo: elektronų sustemoje dalelei, esančiai tam tikroje būsenoje, negalima priskirti konkreių koordinačių.
\end{itemize}

Abu principus tenkina determinantas:
\begin{equation*}
    \Psi^a(\gamma_1 \gamma_2 | \tau_1, \tau_2) = \frac{1}{\sqrt{2}}
    \det \left( 
    \begin{bmatrix*}[c]
        \psi(\gamma_1 | \tau_1) & \psi(\gamma_1 | \tau_2) \\
        \psi(\gamma_2 | \tau_1) & \psi(\gamma_2 | \tau_2)
    \end{bmatrix*}
    \right) = \frac{1}{\sqrt{2}} [\psi(\gamma_1 | \tau_1) \psi(\gamma_2 | \tau_2) - \psi(\gamma_2 | \tau_1) \psi(\gamma_1 | \tau_2)]
\end{equation*}
čia $\gamma_i \equiv n_i l_i m_i \mu_i$ ir $\tau_i \equiv \theta_i, \phi_i, r_i, \sigma_i$. Funkcija keičia ženklą, sukeičiant elektronų koordinates, tokia funkcija vadinama antisimetrine.

...

Determinantinę funkciją galima transformuoti į surištų

...

Antisimetrinė $N$ elektronų nesurištų momentų banginė funkcija

\begin{equation*}
    \Psi^{a} (\gamma_1 \gamma_2 \dots \gamma_N | \tau_1, \tau_2, \dots \tau_N) = 

    \frac{1}{\sqrt{N!}} \det \left(  \right)
\end{equation*}
...

Surištų momentų $N$ ekvivalentinių elektronų antisimetrinė banginė funkcija yra gaunama iš $N-1$ elektrono antisimetrin4s funkcijos ir vieno elektrono funkcijos sandaugų, sudarant tokias jų tiesines kombinacijas, kurios užtikrintų funkcijos antisimetriškumą prijungiamo elektrono atžvilgiu:

\begin{equation*}
    \Psi( l^N \gamma LSM_LM_S) = \sum_{\bar{\gamma} \bat{L} \bar{S}} (l^N \gamma L S || l^{N-1} \bar{\gamma} )
\end{equation*}

...

Pilnutinio jud4jimo kiekio momento banginės funkcijos
Surišant resultatinius momentus $L$ ir $S$, galima pereiti nuo $\Psi(\gamma LSM_LM_S)$ bazė prie $\Psi(\gamma LSJM)$ bazės:
\begin{equation*}
    \Psi(\gamma LSJM) = \Sum_{M_L M_S}
    \begin{bmatrix*}[c]
        L & S & J \\
        M_L & M_S & M
    \end{bmatrix*}
    \Psi(\gamma LSM_L M_S).
\end{equation*}
Kvantinis skaičius $J$ keičiasi nuo $|L-S|$ ligi $L+S$ per vienetą:

\begin{equation*}
    J = |L-S|, |L-S| + 1, \dots, L+S
\end{equation*}
o jo projekcijos įgyja vertes nuo $-J$ ligi $J$, irgi kisdama per vienetą.


\section{Atomo struktūra}

\subsection*{Elektronų sluoksnis ir jo charakteristika}
Elektrono būseną atome aprašo keturi kvantiniai skaičiai:
\begin{itemize}
    \item Pagrindinis kvantinis skaičius $n$,
    \item Orbitinis kvantinis skaičius $l = 0,1,2, \dots, n+1$
    \item Orbitinio momento projekcija $m=0, \pm 1, \dots , \pm l$,
    \item Sukinio projekcija $\mu = \pm 1/2$
\end{itemize}

Atomo elektronai, kurie turi tuos pačius kvantinius skaičius $n,l$ sudaro elektronų sluoksnį. Jie vadinami ekvivalentiniais elektronais.

Pauli principas: du atomo elektronai negali turėti to pačio kvantinių skaičių rinkinio. Tad sluoksnyje gali būti tik tam tikras nedidelis elektronų skaičius.

Projekcija $m$ įgyja $2l+1$ verčių, o $\mu$ - dvi, taigi sluoksnyje $nl^{N}$ gali būti ne daugiau kaip $4l+2$ elektronų. Užpildytas elektronų sluoksnis $l^{4l+2}$. Atviras sluoksnis $l^{N}$ $(0 < N < 5l+2)$:
\begin{itemize}
    \item dalinai užpildytas $N < 2l + 1$
    \item beveik užpildytas $N > 2l + 1$
    \item pusiau užpildytas $N = 2l + 1$
\end{itemize}
Taigi $s^N$ sluoksnyje gali būti ne daugiau...

\subsection*{Būsena, lygmuo ir termas}
Būseną apibūdina pilnas daugiaelektronių kvantinių skaičių rinkinys $\gamma LSJM$, $\gamma$ - kai kurių sluoksnių būsenoms apibūdinti reikalingi papildomi kvantiniai skaičiai. Energijos lygmuo (lygmuo) charakterizuojamas kvantiniais skaičiais $\gamma LSJ$. Būsenų skaičius sluoksnyje:
\begin{equation*}
    g(l^N) = 
\end{equation*}

\subsection*{Termų ir lygmenų nustatymas}
Vienas elektronas $l$ turi vienintelį termą ${}^2K$, pvz

\end{document}
